\documentclass[11pt, a4paper]{article}

\usepackage{ctex}
\usepackage{geometry}
\usepackage{titlesec}
\usepackage{enumitem}
\usepackage{setspace}
\usepackage{hyperref}


\title{Script 梗概}
\author{黄浩宇}
\date{\today}

\begin{document}

        \maketitle
        \newpage
        \tableofcontents

        \newpage

        \section{【梗概】}

        \begin{itemize}[leftmargin=2em, itemsep=0.2em]
                \item \textbf{核心冲突}：学校以“提高学生综合素养”为名，把周六放假改为上课，并开设兴趣培训班；表面上人人赞同，实际上不同意见被压制，视察与新闻则进一步把形式主义包装成教育成果。
                \item \textbf{开场}：画面出现鲨鱼血腥的嘴巴和尖牙，并伴随凄厉惨叫声，奠定压抑、荒诞的基调。
                \item \textbf{第1幕}：校长在会议室提出“提高综合素养”“周六上课、开设兴趣班”等要求。全场领导和老师像机器一样整齐鼓掌、整齐发笑；开明老师提出异议，却被校长用荒谬逻辑反问并遭到全场整齐嘲笑。最终校长宣布散会，所有人像机器人一样僵硬走出会场。
                \item \textbf{第2幕}：星期六，教室里挂着“电路设计兴趣小组”的牌子，学生像机器人一样操作电路板。报时器一响，学生整齐通电、亮灯，并异口同声承认“素质提高了”；随后又整齐离开。黑板上字变为“兴趣小组结束”，相机快门响起，学校网站随即发布新闻，宣称该校周末兴趣小组使学生素质普遍提高，成为全市领先高中。
                \item \textbf{第3幕}：教育部部长到校视察，随机采访学生。学生A、学生B都说兴趣班“很无聊，听不懂”，被秘书拉走；学生C则用标准官话称赞活动“非常有意义”“受益匪浅”，部长立刻开怀大笑并鼓励他。随后，本市教育部发布新闻，宣称全市学生素质普遍提高、认可度极高，已成为全省教育先进市。
                \item \textbf{主题指向}：通过机器人般的掌声、步伐、回答和新闻宣传，讽刺形式主义、统一口径与教育评价中的虚假繁荣。
        \end{itemize}

\end{document}